\section*{Overview}

Prefix sums and difference arrays are among the smallest tools in the toolkit, but they repeatedly remove an entire
logarithm or even an entire data structure from the solution. They are often the right answer before segment trees and
Fenwick trees enter the conversation.

\section*{When to Use It}

Use a prefix sum when:

\begin{itemize}
  \item the query asks for a range sum, count, or any range quantity expressible as a difference of prefixes,
  \item the array is static or only built once,
  \item many queries are asked after a preprocessing step.
\end{itemize}

Use a difference array when:

\begin{itemize}
  \item many updates add the same value on a whole interval,
  \item the final array is only needed after all updates are known,
  \item reconstructing once at the end is acceptable.
\end{itemize}

\section*{Core Idea}

For prefix sums, define
\[
pref[i] = a_1 + a_2 + \cdots + a_i.
\]
Then the sum on \([l, r]\) becomes
\[
pref[r] - pref[l - 1].
\]

For difference arrays, store changes instead of values:
\[
diff[i] = a_i - a_{i-1}.
\]
To add \(x\) on \([l, r]\), do
\[
diff[l] += x, \qquad diff[r + 1] -= x.
\]
One final prefix pass reconstructs the actual array.

\section*{Key Insight}

Both techniques work by moving the work to boundaries.

\begin{itemize}
  \item Prefix sums answer an interval by looking only at its two ends.
  \item Difference arrays mark only where an update starts and where it stops.
\end{itemize}

This boundary view is the reason they are so common in greedy checks, offline processing, and line-sweep style tasks.

\section*{Operations / Main Technique}

\paragraph{Prefix-sum query.}
Build once in \(O(n)\), answer each range sum in \(O(1)\).

\paragraph{Difference-array update.}
Apply each range add in \(O(1)\), rebuild in \(O(n)\) at the end.

\paragraph{Worked example.}
If I need to add \(+3\) on \([2, 5]\) and \(+4\) on \([4, 6]\), the difference array marks only four positions:
\[
diff[2] += 3, \; diff[6] -= 3, \; diff[4] += 4, \; diff[7] -= 4.
\]
The prefix of \(\texttt{diff}\) then recreates the final values.

\section*{Correctness Intuition}

For prefix sums, each interval sum is the full prefix to \(r\) minus the prefix before \(l\). Everything outside the
interval cancels.

For difference arrays, a range update raises the running prefix by \(x\) starting at \(l\), and removes that extra
contribution after \(r\). So every position inside the range sees \(x\), and every position outside does not.

\section*{Complexity Analysis}

\begin{itemize}
  \item build prefix sums: \(O(n)\),
  \item range sum query: \(O(1)\),
  \item one difference-array range update: \(O(1)\),
  \item reconstruct final array from differences: \(O(n)\),
  \item memory: \(O(n)\).
\end{itemize}

\section*{Implementation}

The code includes one small helper for static prefix sums and one helper for offline range additions with a difference
array. Both are 0-indexed externally, which avoids repeated shifts in typical vector-based code.

\section*{Common Pitfalls}

\begin{itemize}
  \item Forgetting the extra slot needed for \(\texttt{diff[r + 1]}\).
  \item Mixing inclusive and half-open intervals.
  \item Reaching for prefix sums when updates are online and frequent.
  \item Using \texttt{int} for sums that should be \texttt{long long}.
  \item Forgetting that prefix sums work for more than sums: counts, XOR, and weighted prefixes show up often too.
\end{itemize}

\section*{Variants / Extensions}

\begin{itemize}
  \item Two-dimensional prefix sums for submatrix queries.
  \item Two-dimensional difference arrays for rectangle updates.
  \item Prefix sums over frequencies after coordinate compression.
  \item Weighted prefix sums for fast range cost formulas.
\end{itemize}

\section*{Practice Problems}

\begin{itemize}
  \item Many static range-sum queries on one array.
  \item Offline interval additions followed by one final scan.
  \item Counting active segments at every coordinate after compression.
  \item Binary-search checkers that need fast range totals.
\end{itemize}

\section*{References}

\begin{itemize}
  \item \href{https://usaco.guide/silver/prefix-sums}{USACO Guide: Prefix Sums}.
  \item \href{https://cp-algorithms.com/data_structures/fenwick.html}{cp-algorithms: Fenwick Tree} for the online extension of the same prefix idea.
  \item \href{https://codeforces.com/catalog}{Codeforces Catalog}.
\end{itemize}
